\documentclass[12pt, a4paper]{article}

\usepackage[margin=1.8cm,top=2cm,bottom=2cm]{geometry}
\usepackage{amsmath,amssymb,amsfonts}
\usepackage{graphicx}
\usepackage[colorlinks=true,linkcolor=blue,citecolor=blue,urlcolor=blue]{hyperref}
\usepackage{bm,xcolor}
\usepackage[numbers,sort&compress]{natbib}
\usepackage{float}

\newcommand{\Geff}{G_{\mathrm{eff}}}
\newcommand{\meff}{m_{\mathrm{eff}}}
\newcommand{\heff}{\hbar_{\mathrm{eff}}}
\newcommand{\rhob}{\bar{\rho}}
\newcommand{\lP}{\ell_{\mathrm{P}}}

\begin{document}

\title{The Quantum Hall Effect as Topological
Phase-Locking in the Logarithmic Superfluid Vacuum}

\author{B. B. Kulangiev\\\small{Haifa, Israel}}
\date{April 2026}
\maketitle

\begin{abstract}
\noindent
We show that the integer quantum Hall effect is a
natural consequence of the logarithmic superfluid
vacuum framework. The electron, modeled as a
vortex-Gausson soliton with winding number $n = -1$,
couples to an external magnetic field through the
gauge-invariant phase velocity
$u_\mu = (\partial_\mu S - eA_\mu)/\meff$. The
applied field $\mathbf{B}$ imposes a background
vorticity $\nabla\times\mathbf{A}$ on the condensate,
to which the electron vortex phase-locks. We derive
four results: (i)~the von Klitzing constant
$R_K = h/e^2$ emerges as the impedance of a single
quantized circulation in the condensate;
(ii)~Landau levels are the discrete vortex-lattice
modes of the electron-condensate coupling, with
spacing $\hbar\omega_c = \hbar eB/\meff$;
(iii)~the quantized Hall conductance
$\sigma_{xy} = \nu e^2/h$ follows from the
topological invariance of the total winding number
$\nu$ of filled vortex modes; and (iv)~the thermal
breakdown scale $k_B T^* \sim \hbar\omega_c$
is recovered as the energy required to break the
phase-locking between the electron vortex and the
background field. All four results reproduce the
standard theory exactly, with no new free parameters.
The framework provides a mechanical interpretation
of topological protection: backscattering is
suppressed because it requires changing the winding
number of a vortex defect, which is a topological
invariant of the condensate.
\end{abstract}


% ====================================================================
\section{Introduction}
\label{sec:intro}
% ====================================================================

The integer quantum Hall effect
(IQHE)~\cite{klitzing_1980} is one of the most
precisely verified phenomena in physics: the Hall
resistance of a two-dimensional electron gas in a
strong perpendicular magnetic field is quantized to
$R_{xy} = h/(\nu e^2)$ with a precision exceeding
one part in $10^9$, where $\nu$ is an integer. The
standard explanation invokes Landau quantization of
the orbital motion and the topological protection of
edge states via the Chern number~\cite{thouless_1982}.

While mathematically complete, the standard
derivation treats the quantization as an abstract
topological property of the single-particle
Hamiltonian in a magnetic field. No mechanical
explanation is offered for \emph{why} backscattering
is forbidden or \emph{what physical structure}
carries the topological invariant.

In the logarithmic superfluid vacuum
framework~\cite{kulangiev_gravity,kulangiev_geff},
the vacuum is a physical condensate described by a
complex order parameter
$\psi = \sqrt{\rho}\,e^{iS/\heff}$, and
fundamental particles are topological defects
(vortex-Gausson solitons) of this
condensate~\cite{kulangiev_planck,kulangiev_inertia}.
In a companion paper~\cite{kulangiev_maxwell}, we
showed that gauging the $U(1)$ phase symmetry of
the condensate produces Maxwell's equations, with
electric charge quantized by the vortex winding
number: $Q = ne$.

In this paper, we show that placing a charged
vortex-soliton (the electron) in an external
magnetic field (background condensate vorticity)
naturally produces all the key features of the IQHE.
The derivation uses only the gauge-invariant
Madelung velocity from Paper~8 and standard
superfluid vortex dynamics. No new assumptions or
parameters are introduced.


% ====================================================================
\section{The Electron as a Vortex in a Background
Field}
\label{sec:setup}
% ====================================================================

\subsection{The gauge-invariant velocity}

From the gauged logarithmic Klein-Gordon
equation~\cite{kulangiev_maxwell}, the Madelung
decomposition yields the gauge-invariant
four-velocity:
\begin{equation}
    u_\mu = \frac{1}{\meff}
    (\partial_\mu S - eA_\mu)\,.
\label{eq:u_gauge}
\end{equation}
The spatial part gives the condensate flow velocity
around the electron:
\begin{equation}
    \mathbf{v} = \frac{1}{\meff}
    (\nabla S - e\mathbf{A})\,.
\label{eq:v_gauge}
\end{equation}
In the absence of an external field
($\mathbf{A} = 0$), the electron vortex with winding
number $n$ has circulation:
\begin{equation}
    \oint \mathbf{v}\cdot d\boldsymbol{\ell}
    = \frac{2\pi n\heff}{\meff}\,.
\label{eq:circulation}
\end{equation}
This is the Onsager-Feynman quantization condition
for a superfluid vortex~\cite{feynman_1955}.

\subsection{The magnetic field as background vorticity}

An external magnetic field $\mathbf{B} = \nabla
\times\mathbf{A}$ imposes a background vorticity on
the condensate. From Eq.~(\ref{eq:v_gauge}), the
vorticity of the condensate flow is:
\begin{equation}
    \boldsymbol{\omega} = \nabla\times\mathbf{v}
    = \frac{1}{\meff}(\nabla\times\nabla S
    - e\nabla\times\mathbf{A})
    = -\frac{e\mathbf{B}}{\meff}
    + \frac{2\pi n\heff}{\meff}
    \delta^2(\mathbf{x})\,\hat{\mathbf{z}}\,,
\label{eq:vorticity}
\end{equation}
where the first term is the smooth background
vorticity from $\mathbf{B}$ and the second is the
singular vorticity of the electron itself (a delta
function at the vortex core). The total vorticity is
the sum of the applied field and the intrinsic
electron circulation.


% ====================================================================
\section{The Von Klitzing Constant as Circulation
Impedance}
\label{sec:RK}
% ====================================================================

In the condensate, each vortex carries one quantum
of circulation $\kappa = h/\meff$. The ``impedance''
of a single circulation quantum---the ratio of the
phase-gradient driving force to the resulting
current---is determined by the charge-circulation
coupling.

The magnetic flux quantum threading one vortex
cell is:
\begin{equation}
    \Phi_0 = \frac{h}{e}\,,
\label{eq:Phi0}
\end{equation}
which follows from the single-valuedness of $\psi$
around a loop enclosing one flux quantum
(the Aharonov-Bohm
condition~\cite{aharonov_1959}). The current
carried by one filled Landau level per unit voltage
drop is $I = e\omega_c/(2\pi) = e^2 B/(2\pi\meff)$.
The voltage per flux quantum is
$V = \Phi_0 \omega_c/(2\pi)$. Their ratio gives:
\begin{equation}
    \boxed{\;R_K = \frac{V}{I} = \frac{h}{e^2}
    \approx 25\,812.807\;\Omega\;}
\label{eq:RK}
\end{equation}
In the condensate picture, $R_K$ is the
\emph{impedance of one circulation quantum coupled
to one charge quantum}. It is the natural resistance
unit of a quantized vortex carrying electric charge
$e$. The von Klitzing constant is to the charged
vortex what the acoustic impedance $Z_0 = \mu_0 c$
is to the bulk condensate~\cite{kulangiev_impedance}.

The connection to the vacuum impedance is:
\begin{equation}
    Z_0 = 2\alpha R_K\,,
\label{eq:Z0_RK}
\end{equation}
where $\alpha = e^2/(4\pi\epsilon_0\hbar c)
\approx 1/137$ is the fine structure constant. In
the condensate framework, $\alpha$ is the coupling
efficiency between the vortex circulation and the
bulk medium~\cite{kulangiev_impedance}: the fraction
of the ideal vortex impedance $R_K$ that appears as
bulk wave impedance $Z_0$.


% ====================================================================
\section{Landau Levels as Vortex-Lattice Modes}
\label{sec:landau}
% ====================================================================

\subsection{The cyclotron frequency}

An electron vortex in a uniform background vorticity
$\boldsymbol{\omega}_B = -e\mathbf{B}/\meff$
orbits at the cyclotron frequency. From the
gauge-invariant equation of
motion~\cite{kulangiev_inertia}, the electron
(mass $m = E_v/c^2$, charge $-e$) in a uniform
$\mathbf{B} = B\hat{\mathbf{z}}$ executes circular
orbits with angular frequency:
\begin{equation}
    \omega_c = \frac{eB}{\meff}\,.
\label{eq:omega_c}
\end{equation}
This is the standard cyclotron frequency, here
interpreted as the angular velocity of the electron
vortex orbiting in the background vorticity field.

\subsection{The discrete spectrum}

In the condensate picture, the electron vortex
phase-locks to the background vorticity. The allowed
states are those where the total phase accumulated
by the electron around its orbit is an integer
multiple of $2\pi$:
\begin{equation}
    \oint (\nabla S - e\mathbf{A})\cdot
    d\boldsymbol{\ell} = 2\pi\heff\,n\,,
    \qquad n = 0, 1, 2, \ldots
\label{eq:phase_quantization}
\end{equation}
This is the Bohr-Sommerfeld condition applied to the
gauge-invariant phase. For circular orbits in a
uniform field, this gives equally spaced energy
levels:
\begin{equation}
    \boxed{\;E_n = \hbar\omega_c
    \left(n + \frac{1}{2}\right)\;}
\label{eq:landau}
\end{equation}
These are the Landau levels. In the condensate
framework, each level corresponds to a mode of the
vortex-lattice coupling: the $n$-th level has $n$
additional circulation quanta wound around the
guiding center. The half-integer zero-point energy
$\hbar\omega_c/2$ arises from the irreducible
circulation of the electron vortex itself
($n = -1$), which persists even in the lowest
orbital mode.

\subsection{Degeneracy and flux quantization}

The degeneracy of each Landau level equals the
number of flux quanta threading the sample:
\begin{equation}
    N_\Phi = \frac{BA}{\Phi_0}
    = \frac{eBA}{h}\,,
\label{eq:degeneracy}
\end{equation}
where $A$ is the sample area. In the condensate
picture, this is the number of vortex cells that
tile the sample when the background vorticity is
$\omega_B = eB/\meff$. Each cell has area
$a_v = \Phi_0/B = h/(eB)$, and the sample holds
$A/a_v = eBA/h$ cells. The electron can be
localized in any one of these cells, giving the
degeneracy.


% ====================================================================
\section{Quantized Hall Conductance from Winding
Number Conservation}
\label{sec:hall}
% ====================================================================

\subsection{The filling factor}

The filling factor $\nu$ is the ratio of electron
density $n_e$ to flux quantum density $n_\Phi$:
\begin{equation}
    \nu = \frac{n_e}{n_\Phi}
    = \frac{n_e h}{eB}\,.
\label{eq:filling}
\end{equation}
When $\nu$ is an integer, exactly $\nu$ Landau levels
are completely filled. Each filled level contributes
$N_\Phi$ states, all carrying the same topological
winding number.

\subsection{The Hall conductance}

The current response to a transverse electric field
$E_x$ in a filled Landau level is:
\begin{equation}
    j_y = \sigma_{xy} E_x\,,
\label{eq:hall_current}
\end{equation}
where $\sigma_{xy}$ is the Hall conductance. Each
filled vortex-lattice mode contributes one quantum
of conductance $e^2/h$, because each circulation
quantum transports one charge quantum per flux
quantum of applied field. For $\nu$ filled levels:
\begin{equation}
    \boxed{\;\sigma_{xy} = \nu\,\frac{e^2}{h}\;}
\label{eq:sigma_xy}
\end{equation}
The Hall resistance is:
\begin{equation}
    R_{xy} = \frac{1}{\sigma_{xy}}
    = \frac{h}{\nu e^2}
    = \frac{R_K}{\nu}\,.
\label{eq:R_xy}
\end{equation}

\subsection{Topological protection}

In the condensate framework, the quantization of
$\sigma_{xy}$ is exact because $\nu$ is a
\emph{winding number}---the total number of filled
circulation modes of the vortex lattice. Winding
numbers are integers by topological necessity:
you cannot have half a vortex.

Backscattering (the process that destroys
quantization in ordinary conductors) requires
transferring an electron from one edge to the
opposite edge, which means moving a vortex across
the entire sample. This requires unwinding and
rewinding the vortex topology across all the
intervening flux cells---a process forbidden unless
sufficient energy is supplied to break the
phase-locking.

This provides the mechanical interpretation of
topological protection that is absent in the
abstract Chern-number formulation: the Hall
quantization is robust because vortex winding
numbers are conserved quantities of the condensate
dynamics, not because of an abstract mathematical
theorem.


% ====================================================================
\section{The Thermal Breakdown Scale}
\label{sec:thermal}
% ====================================================================

The phase-locking between the electron vortex and
the background vorticity is maintained as long as
thermal fluctuations cannot supply enough energy to
promote an electron to the next Landau level. The
critical temperature is:
\begin{equation}
    \boxed{\;k_B T^* \sim \hbar\omega_c
    = \frac{\hbar eB}{\meff}\;}
\label{eq:T_star}
\end{equation}
For $B = 10$~T and $\meff = m_e$:
\begin{equation}
    T^* = \frac{\hbar eB}{m_e k_B}
    \approx \frac{1.055 \times 10^{-34}
    \times 1.602 \times 10^{-19} \times 10}
    {9.109 \times 10^{-31}
    \times 1.381 \times 10^{-23}}
    \approx 13\;\text{K}\,.
\label{eq:T_numerical}
\end{equation}
This is consistent with the experimental observation
that the IQHE requires cryogenic temperatures
($T \lesssim 4$~K for clean samples at $B \sim
10$~T)~\cite{klitzing_1980}.

In the condensate language: at $T > T^*$, thermal
phonons in the condensate have sufficient energy to
break the topological phase-locking between the
electron vortex and the background vorticity. The
vortex ``unlocks'' from the lattice mode and scatters
freely, destroying the quantization.

This result confirms that the framework reproduces
the correct energy scales without new parameters.
The thermal wall that prevents room-temperature
topological protection in the IQHE is:
\begin{equation}
    \frac{k_B T_{\mathrm{room}}}{\hbar\omega_c}
    = \frac{0.025\;\text{eV}}{1.2 \times
    10^{-4}\;\text{eV}} \approx 210\,,
\label{eq:thermal_ratio}
\end{equation}
for $B = 1$~T. Room-temperature operation would
require either $B \sim 200$~T (beyond current
technology) or a condensate with an effective mass
$\meff$ two orders of magnitude lighter than $m_e$.
The framework does not provide a mechanism for either.


% ====================================================================
\section{Discussion}
\label{sec:discussion}
% ====================================================================

\subsection{What is derived and what is reproduced}

This paper derives no new physical predictions. Every
result---$R_K = h/e^2$, $E_n = \hbar\omega_c(n+1/2)$,
$\sigma_{xy} = \nu e^2/h$, and
$T^* \sim \hbar eB/(\meff k_B)$---is standard
condensed matter physics~\cite{klitzing_1980,
thouless_1982}. What is new is the \emph{derivation
pathway}: all four results follow from the single
premise that the electron is a vortex defect in the
vacuum condensate, coupled to the background
vorticity through the gauge-invariant Madelung
velocity~(\ref{eq:u_gauge}).

The value of this exercise is threefold:

(i) It demonstrates that the logarithmic superfluid
vacuum framework, originally constructed to derive
emergent gravity~\cite{kulangiev_gravity,
kulangiev_geff}, is consistent with precision
condensed matter physics at the $10^{-9}$ level.

(ii) It provides a mechanical interpretation of
topological protection: vortex winding numbers are
conserved because the condensate is superfluid
(irrotational except at defects), not because of an
abstract Chern number.

(iii) It identifies the thermal breakdown scale
$T^*$ as the vortex phase-unlocking energy, making
explicit why the IQHE requires cryogenic conditions
and why no simple modification of the framework can
bypass this requirement.

\subsection{Relation to the fractional quantum Hall
effect}

The fractional quantum Hall effect (FQHE), observed
at filling factors $\nu = p/q$ with $q$
odd~\cite{laughlin_1983}, involves strong
electron-electron correlations. In the condensate
framework, the FQHE would correspond to
\emph{composite vortex} states: bound states of
$q$ electron vortices sharing $p$ circulation
quanta. The Laughlin wavefunction
$\prod_{i<j}(z_i - z_j)^q\,e^{-|z|^2/4}$ has the
structure of a product of vortex separations raised
to the $q$-th power, consistent with $q$-fold
winding. A detailed derivation of the FQHE from the
logarithmic condensate is beyond the scope of this
paper but is a natural extension.

\subsection{Relation to the vacuum impedance}

The hierarchy of impedance scales in the framework
is:
\begin{equation}
    R_K = \frac{h}{e^2} \approx 25.8\;\text{k}\Omega
    \quad\xrightarrow{\times 2\alpha}\quad
    Z_0 = \mu_0 c \approx 377\;\Omega\,.
\label{eq:hierarchy}
\end{equation}
The von Klitzing constant is the impedance of one
vortex circulation quantum. The vacuum impedance is
the impedance of the bulk condensate, reduced by the
coupling factor $2\alpha \approx 1/68.5$. This
hierarchy---from the microscopic (vortex) to the
macroscopic (bulk)---is the electromagnetic analogue
of the gravitational Feynman-Onsager relation, where
microscopic Bondi flows ($r^{-2}$) produce the
macroscopic PG flow
($r^{-1/2}$)~\cite{kulangiev_geff}. In both cases,
the collective response of the condensate differs
qualitatively from the sum of individual
contributions.

\subsection{Implications for room-temperature
topological electronics}

Equation~(\ref{eq:thermal_ratio}) quantifies the
thermal wall for topological protection: at 1~T,
room-temperature operation requires overcoming a
factor of 210 in energy scale. The framework
identifies three logically possible (but currently
impractical) routes:

(i) Extreme magnetic fields ($B > 200$~T).
Destructive pulsed fields have reached
$\sim 1200$~T~\cite{nakamura_2018}, but sustained
fields of this magnitude are not available.

(ii) Ultralight effective mass. In Dirac
materials (graphene, topological insulators), the
effective mass near the Dirac point approaches
zero, enhancing $\omega_c$ for a given $B$. The
anomalous quantum Hall effect in
graphene~\cite{novoselov_2005} at $T \sim 300$~K
and $B \sim 45$~T exploits precisely this route.

(iii) A new coupling mechanism beyond the standard
$eB/\meff$ cyclotron energy. The logarithmic
condensate does not provide such a mechanism; the
phase-locking energy is set by $\alpha$, $e$,
$\hbar$, and $B$, all of which are fixed.

The framework is therefore honest about its
limitations: it explains the IQHE but does not
transcend it.

\subsection{The soliton rigidity argument and its
failure}

A natural conjecture within the framework is that
the Gausson's internal binding energy
($E_v = \meff c^2 \approx 0.511$~MeV for the
electron) might provide topological protection
far exceeding the Landau gap
$\hbar\omega_c \sim 10^{-4}$~eV. The argument
would be: a thermal phonon at 0.025~eV cannot
deform the Gausson's phase topology, which is
locked by the logarithmic nonlinearity at the
MeV scale, and therefore cannot induce
backscattering.

This argument conflates two distinct processes.
\emph{Destroying} the electron (unwinding the
vortex topology) costs $\meff c^2 \sim$~MeV and
is indeed inaccessible to thermal phonons.
However, \emph{backscattering} the electron
(translating the intact soliton from one guiding
center to another) does not alter the internal
topology. The soliton remains a perfectly rigid
Gausson throughout; only its center-of-mass
coordinate $\mathbf{X}_{\mathrm{cm}}$ changes.

The center-of-mass equation of motion for any
localized wavepacket in a uniform magnetic field
is~\cite{kohn_1961}:
\begin{equation}
    \meff\,\ddot{\mathbf{X}}_{\mathrm{cm}}
    = e\mathbf{E}
    + e\,\dot{\mathbf{X}}_{\mathrm{cm}}
    \times\mathbf{B}\,.
\label{eq:kohn}
\end{equation}
This is Kohn's theorem: the center-of-mass
dynamics of a charged system in a uniform field
is identical to that of a point particle,
regardless of the internal structure, binding
energy, or nonlinearity of the wavepacket. The
cyclotron frequency is $\omega_c = eB/\meff$ and
the Landau gap is $\hbar\omega_c$, independent of
whether the electron is a point charge, a Gaussian
wavepacket, or a logarithmic Gausson soliton.

The physical picture is clear: the Gausson's
logarithmic binding energy protects its
\emph{shape} (internal modes cost $\sim \meff c^2$
to excite), but not its \emph{position} (the
center-of-mass translation is a zero-energy
Goldstone mode of the internal Hamiltonian). A
thermal phonon cannot shatter the soliton, but it
can nudge it sideways into an adjacent vortex cell.
The nudging energy is the Landau gap
$\hbar\omega_c$, not the binding energy $\meff c^2$.

A violation of Kohn's theorem would require a
nonlinear coupling between the center-of-mass
motion and the local condensate density at the
scale of the magnetic length
$\ell_B = \sqrt{\hbar/(eB)} \sim 10^{-8}$~m.
However, the PG density correction is
$\epsilon_1 = 0$ (Sec.~\ref{sec:landau}), so the
condensate density is uniform at all relevant
scales. The nonlinear coupling vanishes, and
Kohn's theorem holds exactly.

We therefore conclude that the logarithmic
nonlinearity does not provide a mechanism for
room-temperature topological protection beyond
the standard Landau gap. Any proposed device
exploiting condensate-mediated phase-locking
(e.g., a ``Vacuum-Tension Transistor'') must
confront the same thermal wall as the conventional
quantum Hall effect:
$k_B T_{\mathrm{room}}/\hbar\omega_c \approx 210$
at 1~T.

\subsection{Deterministic ontology and quantum mechanics}

The framework provides a deterministic ontology for
quantum mechanics analogous to de Broglie-Bohm pilot
wave theory~\cite{bohm_1952}, with the logarithmic
condensate serving as the physical pilot wave medium.
The order parameter $\psi(\mathbf{x},t)$ evolves
deterministically via the logarithmic Klein-Gordon
equation; particles are Gausson solitons with
definite positions and velocities at all times. Quantum
superposition corresponds to a definite condensate
field configuration $\psi = \psi_A + \psi_B$ that is
a coherent sum of two localized perturbations---both
branches coexist in the real fluid, not in an abstract
Hilbert space. Measurement selects a branch through
decoherence: coupling to macroscopic degrees of
freedom distributes the cross-terms across
$\sim 10^{23}$ environmental modes, rendering
interference unobservable without requiring a
collapse postulate. The Born rule
($P \propto |\psi|^2$) is recovered because
$|\psi|^2 = \rho$ is the physical condensate density,
and soliton configurations in thermal contact with
the condensate relax to quantum equilibrium, as shown
by Valentini~\cite{valentini_1991} for general
pilot wave theories. This interpretation does not
modify any prediction of standard quantum mechanics;
it provides the mechanical substrate that the
Copenhagen interpretation leaves unspecified.

% ====================================================================
\section{Conclusion}
\label{sec:conclusion}
% ====================================================================

We have shown that the integer quantum Hall effect
follows naturally from the logarithmic superfluid
vacuum framework:

(1) The von Klitzing constant $R_K = h/e^2$ is the
impedance of one quantized circulation in the
condensate, linked to the vacuum impedance by
$Z_0 = 2\alpha R_K$.

(2) Landau levels $E_n = \hbar\omega_c(n + 1/2)$
are the discrete vortex-lattice modes of the
electron-condensate coupling, with the zero-point
energy arising from the irreducible electron vortex
circulation.

(3) The quantized Hall conductance
$\sigma_{xy} = \nu e^2/h$ follows from the
topological invariance of the winding number $\nu$
of filled vortex-lattice modes.

(4) The thermal breakdown scale
$k_B T^* \sim \hbar eB/\meff \approx 13$~K at
10~T reproduces the observed cryogenic requirement,
confirming that the framework introduces no new
energy scales.

No new predictions are made. The contribution is the
unification of a precision condensed matter
phenomenon ($10^{-9}$ accuracy) with the emergent
gravity program, using the same condensate, the same
gauge-invariant velocity, and the same topological
vortex structure that produces Newton's law, Maxwell's
equations, and the Planck
scale~\cite{kulangiev_gravity,kulangiev_geff,
kulangiev_maxwell,kulangiev_planck}.


% ====================================================================
\section*{Acknowledgments}
% ====================================================================

The author thanks K.~G.~Zloshchastiev for
foundational work on the logarithmic superfluid
vacuum. Claude (Anthropic, claude-opus-4-6) and Gemini (Google, Gemini 3.1 Pro) were used as editorial and computational tools during manuscript preparation. All scientific content, equations, and
interpretations are the author's own.


\begin{thebibliography}{20}

\bibitem{kohn_1961}
W.~Kohn,
Cyclotron resonance and de Haas-van Alphen
oscillations of an interacting electron gas,
Phys.\ Rev.\ \textbf{123}, 1242 (1961).

\bibitem{bohm_1952}
D.~Bohm,
A suggested interpretation of the quantum theory in
terms of ``hidden'' variables,
Phys.\ Rev.\ \textbf{85}, 166 (1952).

\bibitem{valentini_1991}
A.~Valentini,
Signal-locality, uncertainty, and the subquantum
H-theorem,
Phys.\ Lett.\ A \textbf{156}, 5 (1991).

\bibitem{klitzing_1980}
K.~von Klitzing, G.~Dorda, and M.~Pepper,
New method for high-accuracy determination of the
fine-structure constant based on quantized Hall
resistance,
Phys.\ Rev.\ Lett.\ \textbf{45}, 494 (1980).

\bibitem{thouless_1982}
D.~J.~Thouless, M.~Kohmoto, M.~P.~Nightingale,
and M.~den Nijs,
Quantized Hall conductance in a two-dimensional
periodic potential,
Phys.\ Rev.\ Lett.\ \textbf{49}, 405 (1982).

\bibitem{laughlin_1983}
R.~B.~Laughlin,
Anomalous quantum Hall effect: An incompressible
quantum fluid with fractionally charged excitations,
Phys.\ Rev.\ Lett.\ \textbf{50}, 1395 (1983).

\bibitem{novoselov_2005}
K.~S.~Novoselov \emph{et al.},
Two-dimensional gas of massless Dirac fermions in
graphene,
Nature \textbf{438}, 197 (2005).

\bibitem{nakamura_2018}
D.~Nakamura \emph{et al.},
Record indoor magnetic field of 1200~T generated by
electromagnetic flux-compression,
Rev.\ Sci.\ Instrum.\ \textbf{89}, 095106 (2018).

\bibitem{aharonov_1959}
Y.~Aharonov and D.~Bohm,
Significance of electromagnetic potentials in the
quantum theory,
Phys.\ Rev.\ \textbf{115}, 485 (1959).

\bibitem{feynman_1955}
R.~P.~Feynman,
Application of quantum mechanics to liquid helium,
in \emph{Progress in Low Temperature Physics},
Vol.~1, ed.\ C.~J.~Gorter (North-Holland, 1955),
pp.~17--53.

\bibitem{kulangiev_gravity}
B.~B.~Kulangiev,
Conditions for Emergent Gravitational Light Bending from a Logarithmic Superfluid Vacuum,
Preprint (2026).
\url{https://kulangiev.github.io#paper1}

\bibitem{kulangiev_geff}
B.~B.~Kulangiev,
Emergent Newtonian Gravity and the Gravitational Feynman-Onsager Relation in the Logarithmic Superfluid Vacuum,
Preprint (2026).
\url{https://kulangiev.github.io#combined}

\bibitem{kulangiev_planck}
B.~B.~Kulangiev,
Singularity Regularization and the Emergent Planck Scale in the Logarithmic Superfluid Vacuum,
Preprint (2026).
\url{https://kulangiev.github.io#paper5}

\bibitem{kulangiev_maxwell}
B.~B.~Kulangiev,
Emergent Electromagnetism from the Gauged Logarithmic Superfluid Vacuum,
Preprint (2026).
\url{https://kulangiev.github.io#paper8}

\bibitem{kulangiev_inertia}
B.~B.~Kulangiev,
Emergent Inertia, Mass-Energy Equivalence, and the Helmholtz Decomposition of Charge and Spin in the Logarithmic Superfluid Vacuum,
Preprint (2026).
\url{https://kulangiev.github.io#paper-inertia}

\bibitem{kulangiev_impedance}
B.~B.~Kulangiev,
Environment-Dependent Vacuum Impedance: The
$\mu_0$--$\epsilon_0$ Splitting in the Logarithmic Superfluid Vacuum,
Preprint (2026).
\url{https://kulangiev.github.io#paper10}

\end{thebibliography}

\end{document}
